\documentclass[10pt,a4paper]{article}
\usepackage[utf8]{inputenc}
\usepackage[T1]{fontenc}
\usepackage[a4paper,top=20mm,bottom=20mm,left=16mm,right=16mm,headheight=14pt]{geometry}
\usepackage{fancyhdr}
\usepackage{booktabs}
\usepackage{longtable}
\usepackage{array}
\usepackage{enumitem}
\usepackage{amssymb}
\usepackage{xcolor}
\usepackage[hyphens]{url}
\usepackage[colorlinks=true,urlcolor=blue!60!black,linkcolor=blue!60!black]{hyperref}

\setlength{\parindent}{0pt}
\setlength{\parskip}{3pt}
\renewcommand{\arraystretch}{1.28}

\pagestyle{fancy}
\fancyhf{}
\fancyhead[L]{\small LANGUAGE \& SOCIETY}
\fancyhead[R]{\small TRANSCRIPT ANALYSIS WORKSHEET}
\fancyfoot[C]{\small\thepage}
\renewcommand{\headrulewidth}{0.4pt}

\begin{document}

\begin{center}
    {\large\bfseries LANGUAGE \& SOCIETY PROJECT}\\[3pt]
    {\LARGE\bfseries Transcript Analysis Worksheet}\\[3pt]
    {\normalsize Qualitative analysis of political news discourse}
\end{center}
\vspace{2pt}
\noindent\rule{\textwidth}{0.4pt}
\vspace{4pt}


\section*{Creator 5: Republic TV (\textit{The Debate With Arnab Goswami} --- Nationalist / Pro-Establishment English TV News)}
\begin{itemize}[leftmargin=1.5em, itemsep=3pt, topsep=4pt]
  \item \textbf{Analyst:} Baani Solanki (\texttt{2025114016})
  \item \textbf{Channel / Outlet:} Republic TV (\texttt{@RepublicWorld}, Republic Media Network)
  \item \textbf{Ideological / Editorial Positioning:} Hyper-nationalist, pro-BJP/Hindutva English prime-time broadcast television; deploys prosecutorial anchor monologues ("putting the opposition in the dock"), moral panics, and national-security binaries (\texttt{"Nationalists of Bharat"} vs \texttt{"Deep State"} / \texttt{"rabid, bloodthirsty mobs"} / \texttt{"Taliban"} mindset) to insulate the ruling establishment and counter-attack critics.
  \item \textbf{Transcript Files Referenced:}
  \begin{itemize}[leftmargin=1.4em, itemsep=2pt, topsep=2pt]
    \item Event 1: \texttt{transcripts/event1\_ram\_mandir/05\_republic\_tv\_0H-bruoEkaY.txt}
    \item Event 2: \texttt{transcripts/event2\_neet\_protest/05\_republic\_tv\_HTvb6E\_1Qdc.txt}
  \end{itemize}
\end{itemize}
\vspace{4pt}\noindent\rule{\textwidth}{0.4pt}\vspace{4pt}

\section*{Guidelines: Five features to look for in every event}
\begin{longtable}{@{}>{\raggedright\arraybackslash}p{0.17\textwidth} >{\raggedright\arraybackslash}p{0.20\textwidth} >{\raggedright\arraybackslash}p{0.24\textwidth} >{\raggedright\arraybackslash}p{0.32\textwidth}@{}}
\toprule
\textbf{Feature} & \textbf{What to look for} & \textbf{Illustrative example} & \textbf{Questions to answer} \\
\midrule
\endhead
\textbf{1. Naming} & What are people, groups, actions, or events called? & "Students", "protesters", "agitators"; "irregularity", "scam", "controversy". & What exact label is used? Who or what does it refer to? Who uses it? What alternative label could have been used? \\
\addlinespace[4pt]
\textbf{2. Presupposition} & What is treated as already true or taken for granted? & "Why are students creating chaos again?" may presuppose earlier disruption. & What does the wording assume? Is that assumption supported, disputed, or left unexamined? Does it appear in a question or words such as "again", "still", or "finally"? \\
\addlinespace[4pt]
\textbf{3. Agency} & Who is shown doing something? Who is acted upon? Who disappears? & "Police injured protesters" vs. "Protesters were injured." & Who is the actor? Who experiences the action? Is responsibility named, backgrounded, or omitted? \\
\addlinespace[4pt]
\textbf{4. Antagonisation} & Who is positioned against whom? & "Ordinary citizens" vs. "political troublemakers". & Which groups are contrasted? Is one presented as more legitimate, moral, patriotic, dangerous, or disloyal? Is an "us vs. them" division implied? \\
\addlinespace[4pt]
\textbf{5. Intensification / softening} & What is exaggerated, minimised, softened, or euphemised? & "Violent mob" vs. "peaceful demonstrators"; "incident" instead of a direct description. & Which words heighten emotion or seriousness? Which minimise it? Are metaphors, dramatic adjectives, hedges, or euphemisms used? Does the wording match the evidence? \\
\bottomrule
\end{longtable}
\textit{Reminder: A label or implication is not automatically false or manipulative. Examine the surrounding context and evidence.}

\vspace{4pt}\noindent\rule{\textwidth}{0.4pt}\vspace{4pt}

\section*{Event1. Ram Mandir donation controversy}
\textbf{Focus:} How is the controversy framed---as financial accountability, religious sentiment, political motivation, or something else? Be precise about the specific allegation or episode covered.

\begin{itemize}[leftmargin=1.5em, itemsep=3pt, topsep=4pt]
  \item \textbf{Outlet:} Republic TV (\textit{The Debate With Arnab Goswami})
  \item \textbf{Date:} \texttt{2026-06-26} \textit{(the exact day 8 Trust/counting staff were arrested and Champat Rai \& Anil Mishra resigned)}
  \item \textbf{Video title:} \textit{Debate With Arnab: Yogi's Zero Tolerance For Ram Mandir Donation Money Scamsters} (\url{https://youtu.be/0H-bruoEkaY?si=T9EqDVN39KNILYKg})
  \item \textbf{Timestamp(s):} Relevant Event 1 segment \texttt{[00:54--54:06]} (total duration \texttt{54:30}; brief sponsor jingles \texttt{[00:08--00:53]} and \texttt{[54:12--54:30]} excluded per project guidelines).
  \item \textbf{Analyst:} Baani Solanki (\texttt{2025114016})
  \item \textbf{Language:} English (Native English broadcast audio, Roman script; non-AI-dubbed)
\end{itemize}

\subsection*{Event-specific questions (Event 1)}
\begin{enumerate}[leftmargin=1.6em, itemsep=4pt, topsep=4pt]
  \item \textbf{What terms are used: "scam", "alleged irregularity", "corruption", "political controversy", or something else?}
  \begin{itemize}[leftmargin=1.4em, itemsep=2pt, topsep=2pt]
    \item Arnab Goswami concedes the basic offence in bounded terms---\texttt{"theft from the coffers of the Bhavya Raander [Ram Mandir]"} (\texttt{[00:54--01:02]}), \texttt{"irregularities"} (\texttt{[03:42]}), \texttt{"corruption by low-level functionaries"} (\texttt{[03:22--03:33]}), and \texttt{"donation money scamsters"} (\texttt{[Title]})---while explicitly denying systemic scope: \texttt{"there is no religious angle, there is no political angle, there is no government angle in this controversy"} (\texttt{[04:00--04:08]}).
  \end{itemize}
  \item \textbf{How are the temple, trust, donors, and people making allegations described?}
  \begin{itemize}[leftmargin=1.4em, itemsep=2pt, topsep=2pt]
    \item \textbf{Temple \& State Leadership:} Revered as \texttt{"Bhavya Raander"} (\texttt{[01:02]}); CM Yogi Adityanath is hailed as \texttt{"Yogiad who's proved that there is zero tolerance for any form of corruption"} (\texttt{[03:22]}, \texttt{[09:00]}).
    \item \textbf{Trust \& Arrested Staff:} The 8 arrested men (\texttt{Avinash Shukla}, \texttt{Lavkush Mishra}, \texttt{Anukalp Mishra}, \texttt{Manish Kumar Yadav}, \texttt{Karunesh Pandey}, \texttt{Ram Shankar Mishra}, \texttt{Ram Shankar alias Tinnu Yadav}, \texttt{Subhash Srivastava}) are repeatedly ring-fenced as \texttt{"low-level functionaries in the Ramand trust, some cashiers"} (\texttt{[03:22--03:33]}, \texttt{[07:14]}, \texttt{[21:09]}, \texttt{[32:30--32:41]}). General Secretary Champat Rai (who resigned that day, \texttt{[35:17--35:25]}) is defended during the middle 30 minutes as a \texttt{"79year-old prominent leader of the Vishvah Hindu Parishad... a man who's worked tirelessly"} (\texttt{[20:23--21:28]}), though later at \texttt{[26:47--27:05]} and \texttt{[39:35--40:16]} Arnab briefly questions whether Champat Rai was \texttt{"not skilled administratively"} or given \texttt{"an honorable exit"}.
    \item \textbf{Donors:} Devotees who \texttt{"throng the Aayodhya Ram mandr"} (\texttt{[06:02]}) and \texttt{"contributed, fought for 500 years"} (\texttt{[35:52--36:02]}), exemplified in VTR by a devotee proclaiming \texttt{"I will never lose faith in Ram because of some selfish people"} (\texttt{[09:19--09:30]}).
    \item \textbf{People Making Allegations:} Demonized en masse in the opening monologue as \texttt{"People who say let's divide Hindus... People who did not want the Ram mandr itself... People who said Lord Ram did never exist"} (\texttt{[01:23--01:52]}), \texttt{"don't have the asta [faith] to visit the Ram temple"} (\texttt{[04:18]}), and heirs to Mulayam Singh Yadav's \texttt{"GENOCIDE OF CAR SEAKS"} (\texttt{[04:47--05:06]}).
  \end{itemize}
  \item \textbf{Is the issue framed as financial accountability, an attack on Hindu faith, a political dispute, or a matter of public trust?}
  \begin{itemize}[leftmargin=1.4em, itemsep=2pt, topsep=2pt]
    \item Framed overwhelmingly as \textbf{swift BJP/Yogi administrative triumph (\texttt{"Zero Tolerance"}) under opportunistic political attack by anti-Hindu opposition hypocrites}. Arnab explicitly announces his prosecutorial pivot at \texttt{[05:14--05:24]}: \texttt{"I believe action has been taken, but since they have politicized the matter I suggest the best response is to put them in the dock and ask them what their record is on Hindutva"}.
  \end{itemize}
  \item \textbf{Are questions about donations presented as legitimate scrutiny or as an attempt to malign the temple or its supporters?}
  \begin{itemize}[leftmargin=1.4em, itemsep=2pt, topsep=2pt]
    \item Opposition questions are framed as \textbf{hypocritical attempts to malign the temple and divide Hindus} (\texttt{"looking at this opportunity to target the Ramander itself"}, \texttt{[01:23]}; \texttt{"want to fish in troubled waters"}, \texttt{[23:01]}). Only late in the 54-minute show (\texttt{[39:35--40:26]}, \texttt{[53:10--53:19]}) does Arnab concede that \texttt{"phase by phase, week by week... details must come out"} because devotees were hurt.
  \end{itemize}
  \item \textbf{Are allegations attributed to specific people and supported with evidence, or presented as established facts?}
  \begin{itemize}[leftmargin=1.4em, itemsep=2pt, topsep=2pt]
    \item Official SIT findings---\texttt{8} named arrests, \texttt{"over 80 lakh rupees may have been stolen"}, and \texttt{"CCTV footage reportedly showing the accused taking cash at least 70 times"} (\texttt{[06:48--07:34]})---are presented with cautious journalistic hedging (\texttt{may have been}, \texttt{reportedly}). Conversely, whenever opposition panelists allege wider BJP/Trust complicity (e.g., Priyanka Kakkar citing Raju Manwani's \texttt{200 kg} silver bricks \& Champat Rai refusing PM summons at \texttt{[14:46--15:23]}; Major Ashish Chaturvedi alleging BJP leaders' involvement and \texttt{2020} private audit findings at \texttt{[27:34--28:14]}, \texttt{[31:25]}), Arnab aggressively demands courtroom-grade proof on the spot (\texttt{"NAME ONE... VALIDATE IT RIGHT NOW"}, \texttt{[31:43--32:10]}).
  \end{itemize}
  \item \textbf{How are denials and counterarguments represented: explained, dismissed, challenged, or omitted?}
  \begin{itemize}[leftmargin=1.4em, itemsep=2pt, topsep=2pt]
    \item \textbf{State/BJP defenses are amplified as the show's headline thesis}, with BJP's Nalin Kohli and Shantanu Gupta given extended turns to praise \texttt{"breakneck speed"} arrests (\texttt{[12:47]}, \texttt{[48:20--48:48]}, \texttt{[49:37--50:40]}). Meanwhile, \textbf{opposition counterarguments are derailed for over 20 minutes (\texttt{[16:50--35:00]} \& \texttt{[46:18--48:02]})} when AAP spokesperson Priyanka Kakkar makes a hyperbolically phrased demand (\texttt{"Why are you shielding Champat and the bigger chors? Why should they not be publicly executed?"}, \texttt{[16:50]})---which Arnab and Nalin Kohli seize upon as \texttt{"BREAKING NEWS"}, \texttt{"TALIBAN"}, \texttt{"lynch mob attitude"}, and \texttt{"PART OF THE TALIBAN SHOOTING SQUAD"} (\texttt{[17:13--22:35]}, \texttt{[33:17--33:37]}) to drown out her preceding factual questions about missing silver receipts and unmade public SIT reports!
  \end{itemize}
  \item \textbf{Does the presenter distinguish criticism of the trust or its financial practices from criticism of the temple or religion itself?}
  \begin{itemize}[leftmargin=1.4em, itemsep=2pt, topsep=2pt]
    \item \textbf{Rhetorically conflates opposition criticism of the Trust with hatred of the Temple and Hinduism} for most of the broadcast (\texttt{"looking at this opportunity to target the Ramander itself"}, \texttt{[01:23]}; Nalin Kohli equating Priyanka Kakkar's attack on Champat Rai with \texttt{"contempt for Lord Ram"} and \texttt{"EXECUTE THE RAM BHAKTS, EXECUTE THE SANATANIS"} at \texttt{[24:51--25:21]}), while reserving for \textit{himself} and pro-BJP lawyer Ishkaran Bhandari (\texttt{[35:44--36:56]}, \texttt{[39:35--40:16]}) the right to distinguish administrative failure from faith.
  \end{itemize}
  \item \textbf{Who is portrayed as protecting donors' interests, and who as exploiting religious sentiment?}
  \begin{itemize}[leftmargin=1.4em, itemsep=2pt, topsep=2pt]
    \item \textbf{Protecting donors' interests:} UP CM Yogi Adityanath (\texttt{[01:12]}, \texttt{[03:22]}, \texttt{[09:00]}), the Trust itself (claimed to have written the letter requesting action, \texttt{[03:42--03:51]}), and the three-member IAS/IPS/Finance SIT (\texttt{[48:29--48:39]}).
    \item \textbf{Exploiting religious sentiment:} Arvind Kejriwal (mocked repeatedly for his \texttt{"nani told me"} Babri Masjid remark, \texttt{[02:01--02:28]}, \texttt{[10:34--11:10]}, \texttt{[47:24--47:34]}), Akhilesh Yadav, and Rahul Gandhi (\texttt{[02:36--05:14]}).
  \end{itemize}
\end{enumerate}

\subsection*{Apply the five shared features (Event 1)}
\begin{longtable}{@{}>{\raggedright\arraybackslash}p{0.19\textwidth} >{\raggedright\arraybackslash}p{0.38\textwidth} >{\raggedright\arraybackslash}p{0.37\textwidth}@{}}
\toprule
\textbf{Feature} & \textbf{Exact wording + timestamp} & \textbf{What it does in context} \\
\midrule
\endhead
\textbf{1. Naming} & Title: \texttt{"Yogi's Zero Tolerance For Ram Mandir Donation Money Scamsters"} \newline  \texttt{"low-level functionaries in the Ramand trust, some cashiers"} \texttt{[03:22--03:33]} vs \texttt{"79year-old prominent leader of the Vishvah Hindu Parishad"} \texttt{[20:23]} & Minimizes the 8 arrested insiders by rank (\texttt{low-level functionaries, some cashiers}) to sever upward culpability, while venerating Champat Rai via age and lifelong religious service (\texttt{79-year-old prominent leader}). \\
\addlinespace[4pt]
\textbf{2. Presupposition} & \texttt{"what hypocrisy people like Akileles and Rahul Gandhi who don't even want to visit the Ram mantr... are questioning Yogiad who's proved that there is zero tolerance for any form of corruption"} \texttt{[03:04--03:22]} & Factive verb \texttt{"who's proved"} presupposes as indisputable truth that Yogi's government has zero tolerance for corruption, while exclamative \texttt{"what hypocrisy"} presupposes opposition bad faith. \\
\addlinespace[4pt]
\textbf{3. Agency} & \texttt{"there has been theft from the coffers... But it is also true that there has been very strong and effective action... immediate action by the Yogi government, the BJP government"} \texttt{[00:54--01:12]} & Uses an agentless existential construction for the crime (\texttt{there has been theft}---no BJP/Trust agent named), immediately followed by explicit agentive credit (\texttt{by the Yogi government, the BJP government}) for the arrests. \\
\addlinespace[4pt]
\textbf{4. Antagonisation} & \texttt{"People who say let's divide Hindus... People who did not want the Ram mandr itself... GENOCIDE OF CAR SEAKS"} \texttt{[01:23--01:52]}, \texttt{[04:56]} \newline  \texttt{"are you sponsored by the Taliban... PART OF THE TALIBAN SHOOTING SQUAD?"} \texttt{[33:17--33:37]} & Constructs an existential civilizational enemy (\texttt{anti-Hindu dividers}, \texttt{kar-sevak killers}, \texttt{Taliban shooting squad}) so viewers perceive opposition critics of the scam as enemies of Hindu society itself. \\
\addlinespace[4pt]
\textbf{5. Intensification / softening} & \textbf{Softening (scam size \& scope):} \texttt{"over 80 lak rupees may have been stolen"} \texttt{[07:26]}, \texttt{"there is no political angle, there is no government angle"} \texttt{[04:08]}, \texttt{"protecting him a wee bit"} \texttt{[40:45]} \newline  \textbf{Intensification (panelist slip):} \texttt{"MY GOD, THIS IS TALIBAN... UNFORGIVABLE"} \texttt{[17:20]}, \texttt{[27:05]} & Hedges and minimizes the financial theft (\texttt{may have been}, \texttt{80 lakh}, \texttt{a wee bit}), while blowing up a panelist's rhetorical hyperbole (\texttt{publicly executed}) into a 20-minute \texttt{"Taliban"} moral panic. \\
\bottomrule
\end{longtable}

\subsection*{Short summary (Event 1)}
\begin{itemize}[leftmargin=1.5em, itemsep=3pt, topsep=4pt]
  \item \textbf{What interpretation of the event does this video encourage?} That an unfortunate \texttt{Rs.\,80+ lakh} cash-counting theft by 8 \texttt{"low-level"} employees (\texttt{[03:33]}, \texttt{[07:26]}) has already been decisively resolved by CM Yogi Adityanath's \texttt{"zero tolerance"} arrests and senior resignations (\texttt{[01:12]}, \texttt{[35:25]}), and that the \textit{real} scandal is anti-Ram opposition leaders (\texttt{Kejriwal}, \texttt{Akhilesh}, \texttt{Rahul}) hypocritically exploiting the episode ahead of UP elections (\texttt{[01:23--05:24]}, \texttt{[09:39]}).
  \item \textbf{Whose perspective is foregrounded, challenged, or missing?} \textit{Foregrounded:} Arnab Goswami's pro-Yogi prosecutorial frame, BJP spokesperson Nalin Kohli, Yogi biographer Shantanu Gupta, and pro-BJP lawyer Ishkaran Bhandari. \textit{Challenged:} AAP's Priyanka Kakkar (relentlessly hectored over \texttt{"public execution"}, \texttt{[17:09--35:00]}), SP-leaning Major Ashish Chaturvedi (\texttt{[31:25--32:58]}), and CPI/Left panelist Vivek Srivastava (ambushed about China/Soviet funds at \texttt{[38:12--39:16]}). \textit{Missing:} Unreceipted donors, whistleblower Mahipal Singh, and any scrutiny of the 2021 Ayodhya land deals (cut off at \texttt{[45:22]}).
  \item \textbf{What evidence supports your interpretation? What complicates it?} \textit{Supports:} Verifiable registration of the FIR, arrest of 8 named individuals within days (\texttt{[06:57--07:14]}), CCTV documentation of 70 cash-theft instances (\texttt{[07:26--07:34]}), and same-day resignations of Champat Rai and Anil Mishra (\texttt{[35:25]}). \textit{Complicates:} Even Arnab himself is forced at \texttt{[39:47--40:16]} to admit \texttt{"I refuse to believe that this kind of an organized theft which has happened over many months and years could have happened only by the cash counting staff"}---directly contradicting his own opening monologue (\texttt{[03:22--04:08]})!
\end{itemize}
\vspace{4pt}\noindent\rule{\textwidth}{0.4pt}\vspace{4pt}

\section*{Event2. Students' protest regarding NEET paper leak}
\textbf{Focus:} How does coverage construct the legitimacy of student protest and responsibility for the underlying problem?

\begin{itemize}[leftmargin=1.5em, itemsep=3pt, topsep=4pt]
  \item \textbf{Outlet:} Republic TV (\textit{The Debate / Arnab's Take})
  \item \textbf{Date:} \texttt{2026-07-25} \textit{(the exact evening Education Minister Dharmendra Pradhan resigned and CJP called off the protest)}
  \item \textbf{Video title:} \textit{Rahul The Biggest Loser: Arnab Unravels The Political Agenda Behind CJP's Violent Protest} (\url{https://youtu.be/HTvb6E_1Qdc?si=7aFf_zysFjpvX8af})
  \item \textbf{Timestamp(s):} Relevant Event 2 segment \texttt{[00:00--06:41]} (\texttt{100\%} direct-to-camera monologue by Arnab Goswami; total duration \texttt{06:41}).
  \item \textbf{Analyst:} Baani Solanki (\texttt{2025114016})
  \item \textbf{Language:} English with strategic Hindi insertions (\texttt{khali}, \texttt{saboot}, \texttt{sansad ukhad denge}) (Native English audio, Roman script; non-AI-dubbed)
\end{itemize}

\subsection*{Event-specific questions (Event 2)}
\begin{enumerate}[leftmargin=1.6em, itemsep=4pt, topsep=4pt]
  \item \textbf{Are participants called "students", "candidates", "protesters", "agitators", "mobs", or "political pawns"?}
  \begin{itemize}[leftmargin=1.4em, itemsep=2pt, topsep=2pt]
    \item \textbf{Systematically stripped of student identity and criminalized in extreme terms:} Arnab explicitly declares \texttt{"these were not students, you were fooled... they were not students who did this, the students had evaporated"} (\texttt{[04:58--05:26]}). Instead, he labels them \texttt{"rabid, bloodthirsty mobs"} (\texttt{[01:47]}), \texttt{"the mob"} (\texttt{[01:00]}, \texttt{[03:14]}), \texttt{"the Lynch mob of Abhijit Dipke"} (\texttt{[03:54]}), \texttt{"the goons [goods] of Abhijit Dipke"} (\texttt{[04:19]}, \texttt{[06:21]}), \texttt{"Deep State Dipke"} (\texttt{[00:23]}), and \texttt{"Sonam Wangchuk, very questionable man"} (\texttt{[00:33]}). \textit{(Note: YouTube ASR twice mishears \texttt{"CJP"} as \texttt{"BJP"} at \texttt{[00:07]}, \texttt{[00:51]}, \texttt{[06:01]}---in context Arnab clearly says CJP leadership called off the protest and left the ground.)}
  \end{itemize}
  \item \textbf{Are their stated demands explained, summarised, mocked, questioned, or left out?}
  \begin{itemize}[leftmargin=1.4em, itemsep=2pt, topsep=2pt]
    \item \textbf{Reduced to a brief transactional opening acknowledgement, then twisted into regime conspiracy:} In the first 16 seconds (\texttt{[00:00--00:15]}), Arnab notes \texttt{"Dharmendra Pradhan has resigned. That demand has been met"} plus an assurance against cases for Dipke and Saurabh Das---only to claim for the remaining 6 minutes that the real hidden agenda was \texttt{"streetocracy"} (\texttt{[00:42--00:51]}), \texttt{"sansad ukhad denge, Nepal model"} regime change (\texttt{[05:08--05:16]}), \texttt{"conversion of Hindus"}, and blocking the \texttt{FCRA}, \texttt{delimitation}, and \texttt{one nation, one poll} bills (\texttt{[02:38--02:56]})!
  \end{itemize}
  \item \textbf{Does the presenter focus on the original grievance, disruption, political affiliations, or public inconvenience?}
  \begin{itemize}[leftmargin=1.4em, itemsep=2pt, topsep=2pt]
    \item \textbf{Zero words on the NEET paper leak itself!} The word \texttt{"NEET"} or \texttt{"paper leak"} does not appear \textit{even once} in the spoken transcript (\texttt{[00:00--06:41]}). 100\% of the monologue focuses on (a) mocking Rahul Gandhi (\texttt{[00:00--02:56]}) and (b) inciting permanent nationalist fury over street violence against police/RAF (\texttt{[03:06--06:38]}).
  \end{itemize}
  \item \textbf{Who is represented as responsible: examination authorities, government, students, political parties, or others?}
  \begin{itemize}[leftmargin=1.4em, itemsep=2pt, topsep=2pt]
    \item \textbf{Zero responsibility is attributed to the NTA, Dharmendra Pradhan, or the government.} Instead, blame is placed entirely on \texttt{Rahul Gandhi}, \texttt{Arvind Kejriwal}, \texttt{Yogendra Yadav} (\texttt{[00:33--00:51]}), \texttt{"Deep State Dipke"} (\texttt{[00:23]}, \texttt{[04:19--05:43]}), and \texttt{"questionable"} \texttt{Sonam Wangchuk} (\texttt{[00:33]}).
  \end{itemize}
  \item \textbf{Are questions framed around whether students have legitimate grievances, or why they are causing disruption?}
  \begin{itemize}[leftmargin=1.4em, itemsep=2pt, topsep=2pt]
    \item Framed exclusively around \textbf{opposition frustration and punitive retribution against protest organisers}: \texttt{"what exactly is your problem, Rahul Gandhi? Why do you sound so frustrated?"} (\texttt{[01:00--01:19]}) and \texttt{"Dharmendra Pradhan has resigned, when will Abhijit Dipke be held to account now?"} (\texttt{[05:34]}).
  \end{itemize}
  \item \textbf{Are claims of violence, disorder, or political manipulation attributed to sources and substantiated?}
  \begin{itemize}[leftmargin=1.4em, itemsep=2pt, topsep=2pt]
    \item \textbf{Presented as unattributed, emotive assertions coupled with extreme geopolitical/communal leaps:} While referencing visual clips (\texttt{"Do not forget these pictures"}, \texttt{[03:26]}, \texttt{[04:03]}) and naming injured RAF officer \texttt{Sonia Sherawat} (\texttt{[03:44--04:11]}), Arnab offers \textbf{zero evidence} for \texttt{"Deep State Dipke"} (\texttt{[00:23]}), \texttt{"Nepal model"} overthrow (\texttt{[05:16]}), or the bizarre claim at \texttt{[02:38--02:48]} that protesters/opposition \texttt{"didn't want the FCRA bill to be brought in. They wanted the conversion of Hindus"}!
  \end{itemize}
  \item \textbf{How are police and government actions described compared with protesters' actions? Are similar actions evaluated differently?}
  \begin{itemize}[leftmargin=1.4em, itemsep=2pt, topsep=2pt]
    \item \textbf{Total inversion of victim and perpetrator:} The government is praised five times in a row via anaphora (\texttt{"the government outplayed them"}, \texttt{[02:06--02:56]}), and \texttt{"our police, our RAF, our security personnel"} (\texttt{[01:37--01:47]}, \texttt{[03:36--04:11]}) are portrayed purely as heroic victims of an \texttt{"attempt to lynch"}, while \textbf{100\% of police violence} (lathi-charges, 25+ tear-gas shells, pellet guns, Wangchuk's forced removal, unbadged men) is \textbf{completely erased}. Protesters' actions are characterized as \texttt{"absolute anarchy"}, \texttt{"drunken violence"}, and \texttt{"sexual perversion"} (\texttt{[01:37]}, \texttt{[04:46--04:58]}).
  \end{itemize}
  \item \textbf{Are students allowed to explain their demands in their own words, or does the presenter mainly speak for them?}
  \begin{itemize}[leftmargin=1.4em, itemsep=2pt, topsep=2pt]
    \item \textbf{Zero student voices.} Arnab Goswami speaks uninterrupted for all \texttt{06:41}, explicitly denying that genuine students were even present (\texttt{"the students had evaporated"}, \texttt{[05:26]}).
  \end{itemize}
  \item \textbf{Does the outlet's framing differ between coverage of the NEET paper-leak issue before the protest and coverage once students begin protesting?}
  \begin{itemize}[leftmargin=1.4em, itemsep=2pt, topsep=2pt]
    \item \textbf{Complete erasure of the educational issue once street protests challenge the government:} Even on the day the Union Education Minister resigns over NEET (\texttt{[00:00]}), Arnab never mentions exam leaks, student suicides, or NTA reform---turning a massive government concession (\texttt{Dharmendra Pradhan has resigned}) into a triumphalist \texttt{"Rahul The Biggest Loser"} victory lap paired with a call for street counter-mobilization (\texttt{[06:10--06:38]}).
  \end{itemize}
\end{enumerate}

\subsection*{Apply the five shared features (Event 2)}
\begin{longtable}{@{}>{\raggedright\arraybackslash}p{0.19\textwidth} >{\raggedright\arraybackslash}p{0.38\textwidth} >{\raggedright\arraybackslash}p{0.37\textwidth}@{}}
\toprule
\textbf{Feature} & \textbf{Exact wording + timestamp} & \textbf{What it does in context} \\
\midrule
\endhead
\textbf{1. Naming} & \texttt{"Deep State Dipke"} \texttt{[00:23]}; \texttt{"Sonam Wangchuk, very questionable man"} \texttt{[00:33]}; \texttt{"streetocracy"} \texttt{[00:42]}; \texttt{"rabid, bloodthirsty mobs"} \texttt{[01:47]}; \texttt{"Lynch mob / goons of Abhijit Dipke"} \texttt{[03:54]}, \texttt{[04:19]} & Dehumanizes protesters with animalistic disease imagery (\texttt{rabid, bloodthirsty}), conspiracy alliteration (\texttt{Deep State Dipke}), and anti-democratic coinage (\texttt{streetocracy}), erasing student citizenship. \\
\addlinespace[4pt]
\textbf{2. Presupposition} & \texttt{"Arnab Unravels The Political Agenda Behind CJP's Violent Protest"} \texttt{[Title]} \newline  \texttt{"Why are you, as Gen Z would say, so shook tonight?"} \texttt{[01:19]} \newline  \texttt{"these were not students, you were fooled"} \texttt{[04:58]} & Definite NP \texttt{"The Political Agenda Behind CJP's Violent Protest"} presupposes both hidden political orchestration and inherent protest violence; \texttt{"you were fooled"} presupposes viewer deception if they felt sympathy for students. \\
\addlinespace[4pt]
\textbf{3. Agency} & \texttt{"our police, our RAF, our security personnel were vehemently attacked by rabid, bloodthirsty mobs, comes a resignation"} \texttt{[01:37--01:56]} & Passivizes security forces purely as patients/victims (\texttt{were vehemently attacked}), deletes all police use of force, and nominalizes Pradhan's forced exit as an agentless motion verb (\texttt{comes a resignation}---avoiding \textit{"CJP forced the government to sack Pradhan"}!). \\
\addlinespace[4pt]
\textbf{4. Antagonisation} & \texttt{"when this violence was being carried out by the goons of Abhijit Dipke, our brave soldiers were fighting terrorists IN ANANTNAG... Nationalists, mobilize yourself... This nation is ours."} \texttt{[04:19--04:37]}, \texttt{[06:21--06:38]} & Juxtaposes Jantar Mantar protesters with \texttt{"terrorists in Anantnag"} killing Indian soldiers (compare Operation Sindoor callback at \texttt{[01:19--01:37]}), and ends with an exclusionary majoritarian call to arms (\texttt{Nationalists, mobilize yourself... This nation is ours}). \\
\addlinespace[4pt]
\textbf{5. Intensification / softening} & \textbf{Anaphoric intensification:} 5x \texttt{"the government outplayed them"} \texttt{[02:06--02:56]}; 9x \texttt{"Do not forget / Do not forgive"} \texttt{[03:06--06:21]}; \texttt{"rabid, bloodthirsty"}, \texttt{"drunken violence"}, \texttt{"sexual perversion"} \texttt{[01:47]}, \texttt{[04:46--04:58]} \newline  \textbf{Softening (govt defeat):} \texttt{"comes a resignation"} \texttt{[01:47]} & Liturgical, drumbeat repetition (\texttt{Do not forgive, do not forget}) and visceral disgust triggers (\texttt{rabid}, \texttt{drunken}, \texttt{perversion}) whip up enduring grievance, while \texttt{"comes a resignation"} + \texttt{"outplayed"} reframes ministerial capitulation into strategic genius. \\
\bottomrule
\end{longtable}

\subsection*{Short summary (Event 2)}
\begin{itemize}[leftmargin=1.5em, itemsep=3pt, topsep=4pt]
  \item \textbf{What interpretation of the event does this video encourage?} That the CJP protest had nothing to do with NEET students (\texttt{"the students had evaporated"}, \texttt{[05:26]}), was a \texttt{"Deep State"} / Rahul Gandhi \texttt{"streetocracy"} plot (\texttt{[00:23--00:51]}) to unleash \texttt{"rabid, bloodthirsty mobs"} against police/RAF and block nationalist bills like \texttt{FCRA} (\texttt{[01:47]}, \texttt{[02:38--02:56]}), that the Modi government brilliantly \texttt{"outplayed"} the opposition by ending the protest via Pradhan's resignation (\texttt{[02:06--02:56]}), and that \texttt{"Nationalists of Bharat"} must \texttt{"never forgive and never forget"} Dipke's \texttt{"goons"} (\texttt{[03:06--06:38]}).
  \item \textbf{Whose perspective is foregrounded, challenged, or missing?} \textit{Foregrounded:} Exclusively Arnab Goswami's militant nationalist perspective, aligned with the Modi government and injured RAF personnel (\texttt{Sonia Sherawat}, \texttt{[03:44--04:11]}). \textit{Challenged/Demonized:} Rahul Gandhi, Arvind Kejriwal, Yogendra Yadav, Abhijit Dipke, Sonam Wangchuk, and all Jantar Mantar protesters. \textit{Missing:} \textbf{100\% of student grievances}, NEET paper leaks, student suicides, MEA's own denial of foreign funding (cited by Sarthak \texttt{EM1GeZufO6E} \texttt{[00:40]}), and all police lathi/tear-gas/pellet violence!
  \item \textbf{What evidence supports your interpretation? What complicates it?} \textit{Supports:} Confirmed resignation of Dharmendra Pradhan and emptying of Jantar Mantar (\texttt{[00:00--00:23]}), plus real injuries sustained by police/RAF personnel during the 20 July clashes (\texttt{[03:14--04:11]}, corroborated by Aaj Tak \texttt{JXLIs0WtJAM} \texttt{[18:09]}). \textit{Complicates:} An obvious internal logical contradiction---if protesters were truly \texttt{"Deep State"} \texttt{"rabid, bloodthirsty mobs"} and \texttt{"not students"}, why did the Union Government accept their demands, remove its Education Minister, and promise \textit{not} to prosecute Dipke and Saurabh Das (\texttt{[00:00--00:15]})? Furthermore, MEA officially denied foreign funding (\texttt{EM1GeZufO6E}, \texttt{[00:40--00:48]}), refuting \texttt{"Deep State Dipke"}.
\end{itemize}
\vspace{4pt}\noindent\rule{\textwidth}{0.4pt}\vspace{4pt}

\section*{Common form: Blank excerpt record}
\textit{Copy or print this page whenever you find a passage worth analysing. Keep the exact wording and enough context to interpret it.}


\subsection*{Excerpt Record 1 (Event 1 --- Opening Monologue: Minimizing Theft \& Putting Critics "In the Dock")}
\begin{itemize}[leftmargin=1.5em, itemsep=3pt, topsep=4pt]
  \item \textbf{Outlet:} Republic TV (\textit{The Debate}) | \textbf{Date:} \texttt{2026-06-26}
  \item \textbf{Video title:} \textit{Debate With Arnab: Yogi's Zero Tolerance For Ram Mandir Donation Money Scamsters}
  \item \textbf{Timestamp(s):} \texttt{[00:54--01:52]} \& \texttt{[03:22--05:24]} | \textbf{Analyst:} Baani Solanki | \textbf{Language:} English (Roman script)
  \item \textbf{Exact transcript excerpt:}
    \begin{quote}\small\ttfamily [00:54] Yes, ladies and gentlemen, there has been theft from the coffers of the [01:02] Bhavya Raander [Ram Mandir]. It is true. It is unfortunate. But it is also true that there has been very strong and effective action, [01:12] sackings and arrests and resignations, immediate action by the Yogi government, the BJP government. But it is also true [01:23] that some people are looking at this opportunity to target the Ramander itself. People who say let's divide Hindus... [03:22] Yogiad who's proved that there is zero tolerance for any form of corruption including the corruption by low-level [03:33] functionaries in the Ramand trust some cashiers... [04:08] there is no religious angle there is no political angle there is no government angle in this controversy... [05:14] since they have politicized the matter I suggest the best response is to put them [05:24] in the dock and ask them what their record is on Hindua [Hindutva].\end{quote}
  \item \textbf{Feature(s):} $\boxtimes$ Naming $\boxtimes$ Presupposition $\boxtimes$ Agency $\boxtimes$ Antagonisation $\boxtimes$ Intensification / softening
  \item \textbf{Who is speaking? Is this the outlet's wording or a quotation?} Anchor Arnab Goswami's scripted opening editorial monologue.
  \item \textbf{What does the wording imply, foreground, soften, or leave unstated?} Uses a classic \textbf{\texttt{Yes... But...} concession-reversal structure}: concedes the undeniable theft in two short existential clauses (\texttt{"there has been theft... It is unfortunate"}), confines culprits to \texttt{"low-level functionaries, some cashiers"}, credits \texttt{"immediate action by the Yogi government"}, and then pivots the entire 54-minute agenda to prosecuting \textit{the opposition} (\texttt{"put them in the dock"}).
  \item \textbf{What context or evidence supports or complicates this interpretation?} Supported by the arrest of 8 staffers on 26 June (\texttt{[06:57--07:14]}). Complicated by the fact that allegations had circulated since 2021 (Mahipal Singh, private audit) and 7 June 2026 (Akhilesh Yadav tweet), and that Champat Rai and Anil Mishra also resigned that very day (\texttt{[35:25]})---hardly \texttt{"some cashiers"}.
  \item \textbf{How does another outlet describe the same event or actor?} \textbf{The Wire} (\texttt{7rUL8fGgRJA}, \texttt{[03:29--04:59]}) frames the exact same event as \texttt{"massive theft or dacoity"} exposing the \texttt{"moral degeneration"} of the \texttt{"entire Sangh Parivar"}, whereas \textbf{Republic TV} declares \texttt{"there is no political angle, there is no government angle"} (\texttt{[04:08]}).
  \item \textbf{My interpretation, with reasons:} Explicitly saying \texttt{"the best response is to put them in the dock"} (\texttt{[05:14--05:24]}) lays bare Republic TV's discursive function: transforming accountability \textit{of} the ruling establishment into an ideological trial \textit{of} the opposition.
\end{itemize}

\subsection*{Excerpt Record 2 (Event 1 --- Manufacturing a "Taliban" Moral Panic Over Panelist's Hyperbole)}
\begin{itemize}[leftmargin=1.5em, itemsep=3pt, topsep=4pt]
  \item \textbf{Outlet:} Republic TV (\textit{The Debate}) | \textbf{Date:} \texttt{2026-06-26}
  \item \textbf{Video title:} \textit{Debate With Arnab: Yogi's Zero Tolerance For Ram Mandir Donation Money Scamsters}
  \item \textbf{Timestamp(s):} \texttt{[16:50--17:28]}, \texttt{[20:12--21:28]} \& \texttt{[33:17--33:37]} | \textbf{Analyst:} Baani Solanki | \textbf{Language:} English (Roman script)
  \item \textbf{Exact transcript excerpt:}
    \begin{quote}\small\ttfamily [16:50] [AAP's Priyanka Kakkar:] But the issue today is why are you shielding champ [Champat Rai] and the bigger chores? Why are why should they not be publicly executed?\end{quote}
    \begin{quote}\small\ttfamily [16:59] [Arnab \& BJP's Nalin Kohli:] publicly executed. You are calling for the public execution... [17:13] breaking news for the public execution [17:20] of Champat Rai... MY GOD, THIS IS TALIBAN...\end{quote}
    \begin{quote}\small\ttfamily [20:23] [Arnab Goswami:] the Aam Aadmi Party has called for the public execution of 79year-old prominent leader of the Vishvah Hindu Parishad... [21:09] There are eight people who have been arrested but they are cash counting staff...\end{quote}
    \begin{quote}\small\ttfamily [33:17] [Arnab Goswami:] you've called for the public execution are you sponsored by [33:26] the Taliban is the Taliban sponsoring you is the Taliban your boss? Are you ideologically driven by the Taliban? ... Are you [33:37] PART OF THE TALIBAN SHOOTING SQUAD?\end{quote}
  \item \textbf{Feature(s):} $\boxtimes$ Naming $\boxtimes$ Presupposition $\square$ Agency $\boxtimes$ Antagonisation $\boxtimes$ Intensification / softening
  \item \textbf{Who is speaking? Is this the outlet's wording or a quotation?} AAP spokesperson Priyanka Kakkar (\texttt{[16:50]}), BJP spokesperson Nalin Kohli (\texttt{[17:20--20:12]}), and anchor Arnab Goswami (\texttt{[17:13]}, \texttt{[20:23--21:28]}, \texttt{[33:17--33:37]}).
  \item \textbf{What does the wording imply, foreground, soften, or leave unstated?} Just before \texttt{[16:50]}, Priyanka Kakkar asked five concrete evidentiary questions (\texttt{[14:03--15:23]}: unpublished SIT report, Brij Bhushan's silence, PM summons to Champat Rai, Raju Manwani's \texttt{200 kg} unreceipted silver). The instant she utters the reckless figurative phrase \texttt{"publicly executed"}, Arnab declares \texttt{"breaking news"}, re-names Champat Rai a vulnerable \texttt{"79year-old prominent leader"}, and hurls five consecutive \texttt{"Taliban"} questions (\texttt{[33:17--33:37]})---ensuring none of her substantive questions ever have to be answered!
  \item \textbf{What context or evidence supports or complicates this interpretation?} Calling for \texttt{"public execution"} on national TV \textit{is} genuinely unconstitutional hyperbole (as Ishkaran Bhandari soberly notes at \texttt{[36:33]}: \texttt{"We do not have Taliban style punishments... as per law all of them need to be... punished"}). Yet weaponizing that 3-word slip for 20+ minutes to accuse an Indian political party of being \texttt{"PART OF THE TALIBAN SHOOTING SQUAD"} is textbook diversionary intensification.
  \item \textbf{How does another outlet describe the same event or actor?} \textbf{Dhruv Rathee} (\texttt{oeH0fOSJuX8}, \texttt{[00:17--01:20]}) foregrounds the exact Sindhi community \texttt{210 kg} silver-brick story that Priyanka Kakkar raised at \texttt{[15:04--15:23]}, which Republic TV buried beneath the \texttt{"Taliban"} shouting match.
  \item \textbf{My interpretation, with reasons:} This exchange illustrates conversational floor-control in television debate: seizing on an opponent's tonal overreach to shift the topic from \textit{temple corruption} to \textit{opposition extremism}.
\end{itemize}

\subsection*{Excerpt Record 3 (Event 1 --- Arnab's Surprising Late-Show Pivot Questioning "Honorable Exit")}
\begin{itemize}[leftmargin=1.5em, itemsep=3pt, topsep=4pt]
  \item \textbf{Outlet:} Republic TV (\textit{The Debate}) | \textbf{Date:} \texttt{2026-06-26}
  \item \textbf{Video title:} \textit{Debate With Arnab: Yogi's Zero Tolerance For Ram Mandir Donation Money Scamsters}
  \item \textbf{Timestamp(s):} \texttt{[39:35--40:26]} | \textbf{Analyst:} Baani Solanki | \textbf{Language:} English (Roman script)
  \item \textbf{Exact transcript excerpt:}
    \begin{quote}\small\ttfamily [39:35] [Arnab Goswami to BJP's Nalin Kohli:] So Nalin, are you protecting Champatra? Question one. Allow me to complete [39:47] the highest administrative level of the trust must be blamed, not the lowest administrative level. I refuse to believe that this kind of an [39:57] organized theft which has happened over many months and years could have happened only by the cash counting staff. [40:05] Therefore, I want to know whether you will call him for questioning, whether his name will be included in the FIR or whether you are trying to give [40:16] him an honorable exit by telling him to quietly resign, right, and lay the blame on others.\end{quote}
  \item \textbf{Feature(s):} $\boxtimes$ Naming $\boxtimes$ Presupposition $\boxtimes$ Agency $\square$ Antagonisation $\boxtimes$ Intensification / softening
  \item \textbf{Who is speaking? Is this the outlet's wording or a quotation?} Anchor Arnab Goswami questioning BJP spokesperson Nalin Kohli directly (\texttt{[39:35--40:26]}).
  \item \textbf{What does the wording imply, foreground, soften, or leave unstated?} Remarkably, 40 minutes into the broadcast, Arnab briefly adopts the \textit{exact institutional critique} voiced by opposition/independent critics (\texttt{"I refuse to believe... could have happened only by the cash counting staff"}, \texttt{"honorable exit"}), before softening it seconds later (\texttt{"seen to be protecting him a wee bit"}, \texttt{[40:45]}) and letting Nalin Kohli reassure viewers that \texttt{"no one... will be spared"} (\texttt{[42:08--42:25]}).
  \item \textbf{What context or evidence supports or complicates this interpretation?} Crucial complicating evidence for our analysis! It shows why careful line-by-line transcript reading matters: even a fiercely pro-government anchor felt genuine Hindu viewer anger over Ram Mandir theft (\texttt{"people have been hurt by it"}, \texttt{[53:19]}) and could not sustain a 100\% uncritical whitewash of Champat Rai for the full hour.
  \item \textbf{How does another outlet describe the same event or actor?} Matches \textbf{Dhruv Rathee} (\texttt{oeH0fOSJuX8}, \texttt{[13:05--13:14]}: \texttt{"Is it enough to resign? ... should they also be arrested?"}) and \textbf{Sarthak Goswami} (\texttt{lb7gZnyZH0c}, \texttt{[00:46]}: \texttt{"tinnoo draivar bali ka bakra"}), though Arnab hedges it (\texttt{"a wee bit"}) and pairs it with praise for CM Yogi.
  \item \textbf{My interpretation, with reasons:} On an issue touching Hindu devotees' sacred donations, Republic TV performs "bounded accountability"---pressing BJP (\texttt{"a wee bit"}) at minute 40 to maintain credibility with devout viewers while protecting CM Yogi and PM Modi throughout.
\end{itemize}

\subsection*{Excerpt Record 4 (Event 2 --- Spinning Pradhan's Resignation as "The Government Outplayed Them")}
\begin{itemize}[leftmargin=1.5em, itemsep=3pt, topsep=4pt]
  \item \textbf{Outlet:} Republic TV | \textbf{Date:} \texttt{2026-07-25}
  \item \textbf{Video title:} \textit{Rahul The Biggest Loser: Arnab Unravels The Political Agenda Behind CJP's Violent Protest}
  \item \textbf{Timestamp(s):} \texttt{[00:00--00:51]} \& \texttt{[01:37--02:56]} | \textbf{Analyst:} Baani Solanki | \textbf{Language:} English (Roman script)
  \item \textbf{Exact transcript excerpt:}
    \begin{quote}\small\ttfamily [00:00] The biggest loser is Rahul Gandhi. Let me tell you why, ladies and gentlemen. Dharmendra Pradhan has resigned. That demand has been met. The protest has [00:07] been called off by the CJP [ASR: BJP] leadership... [00:23] Not Dipke, not Deep State Dipke, [00:33] not Sonam Wangchuk, very questionable man, but this man, 57-year-old Rahul Gandhi... because he wants streetocracy...\end{quote}
    \begin{quote}\small\ttfamily [01:37] And now, after there was absolute anarchy on the streets, our police, our RAF, our security personnel were [01:47] vehemently attacked by rabid, bloodthirsty mobs, comes a resignation, the very thing that the students [01:56] demanded. But Rahul Gandhi is not happy with that... [02:06] He has been outplayed by the government. Rahul Gandhi, you've been outplayed... [02:38] They didn't want the FCRA bill to be brought in. They wanted the conversion of Hindus, but the government outplayed them.\end{quote}
  \item \textbf{Feature(s):} $\boxtimes$ Naming $\boxtimes$ Presupposition $\boxtimes$ Agency $\boxtimes$ Antagonisation $\boxtimes$ Intensification / softening
  \item \textbf{Who is speaking? Is this the outlet's wording or a quotation?} Anchor Arnab Goswami's direct-to-camera monologue.
  \item \textbf{What does the wording imply, foreground, soften, or leave unstated?} When a Cabinet Minister is forced to resign after a 35-day street protest, normal political grammar frames the government as yielding. Arnab flips defeat into triumph via three linguistic moves: (1) agentless nominalization \texttt{"comes a resignation"} (\texttt{[01:47]}), (2) shifting the grammatical subject of defeat onto \texttt{"Rahul Gandhi"} (\texttt{"The biggest loser"}, 5x \texttt{"outplayed"}), and (3) inventing sensational communal stakes (\texttt{"They wanted the conversion of Hindus"}, \texttt{[02:38]}).
  \item \textbf{What context or evidence supports or complicates this interpretation?} Supported only by the fact that Pradhan's resignation successfully ended the Jantar Mantar sit-in before the rest of the Parliament session (\texttt{[00:07--00:23]}). Complicated by the glaring contradiction of calling protesters \texttt{"Deep State"} \texttt{"rabid, bloodthirsty mobs"} in the very same sentence that admits \texttt{"comes a resignation, the very thing that the students demanded"} (\texttt{[01:47--01:56]})!
  \item \textbf{How does another outlet describe the same event or actor?} \textbf{Sarthak Goswami} (\texttt{EM1GeZufO6E}, \texttt{[00:08--00:58]}) frames the exact same resignation the next day as \texttt{"CJP won... aakhiri mein eego todna hi pada"} [finally the government's ego had to break], and directly mocks \texttt{"doosra Goswami"} at \texttt{[22:22--23:04]}.
  \item \textbf{My interpretation, with reasons:} This passage is a textbook demonstration of \textbf{spin via Agency \& Naming}: turning a historic student victory over the government into \texttt{"Rahul The Biggest Loser"}.
\end{itemize}

\subsection*{Excerpt Record 5 (Event 2 --- Erasing Students \& Linking Jantar Mantar to Anantnag Terrorists)}
\begin{itemize}[leftmargin=1.5em, itemsep=3pt, topsep=4pt]
  \item \textbf{Outlet:} Republic TV | \textbf{Date:} \texttt{2026-07-25}
  \item \textbf{Video title:} \textit{Rahul The Biggest Loser: Arnab Unravels The Political Agenda Behind CJP's Violent Protest}
  \item \textbf{Timestamp(s):} \texttt{[03:36--05:26]} | \textbf{Analyst:} Baani Solanki | \textbf{Language:} English (Roman script)
  \item \textbf{Exact transcript excerpt:}
    \begin{quote}\small\ttfamily [03:36] Do not forget that there was an attempt to lynch an Rapid Action Force personnel... [03:44] Nationalists do not forget that RAF officer Sonia Sherawat [03:54] she was harassed she was targeted she was beaten by the Lynch mob of Abhijit Dipke... [04:19] when this violence was being carried out by the goons of Abhijit Dipke our brave soldiers were fighting terrorists IN ANANTNAG... [04:46] while you were carrying out this drunken violence a soldier risked and lost his life... [04:58] this sexual perversion shown at Jantar Mantar in the name of students these were not students you were fooled... [05:08] they say sansad ukhad denge Nepal [05:16] model they wanted to bring in India... [05:26] they were not students who did this the students had evaporated.\end{quote}
  \item \textbf{Feature(s):} $\boxtimes$ Naming $\boxtimes$ Presupposition $\boxtimes$ Agency $\boxtimes$ Antagonisation $\boxtimes$ Intensification / softening
  \item \textbf{Who is speaking? Is this the outlet's wording or a quotation?} Arnab Goswami's direct monologue.
  \item \textbf{What does the wording imply, foreground, soften, or leave unstated?} Deploys \textbf{temporal juxtaposition (\texttt{when... while...})} to morally equate domestic protesters in Delhi with Kashmir terrorists killing an Indian soldier in Anantnag (\texttt{[04:19--04:46]}), while erasing genuine NEET examinees via \texttt{"the students had evaporated"} (\texttt{[05:26]}).
  \item \textbf{What context or evidence supports or complicates this interpretation?} Injuries to police/RAF personnel (and peripheral hooliganism near CP noted even by Sarthak Goswami \texttt{EM1GeZufO6E} \texttt{[07:12--07:21]}) had a factual kernel, which Arnab generalizes to the \textit{entire} protest while omitting 100\% of police pellet-gun/lathi injuries to students and doctors (\texttt{Mohak Mangal} \texttt{v-RikSqb0Vw} \texttt{[01:21--02:59]}; \texttt{The Wire} \texttt{pYaTVvHeGoM} \texttt{[10:07--10:38]}).
  \item \textbf{How does another outlet describe the same event or actor?} \textbf{Mohak Mangal} (\texttt{v-RikSqb0Vw}, \texttt{[01:03]}, \texttt{[06:06]}, \texttt{[10:56]}) samples this exact Republic TV rhetoric (\texttt{"brutal violence carried out by the goons of the CJP"}, \texttt{"these are not students"}) and refutes it solely by letting real NEET repeaters, mothers, and teachers speak on camera.
  \item \textbf{My interpretation, with reasons:} Invoking martyred soldiers in Anantnag (\texttt{[04:19--04:46]}) and Operation Sindoor (\texttt{[01:24]}) during an exam-leak story is classic \textbf{national-security displacement}: making any sympathy for student protesters feel like treason against the armed forces.
\end{itemize}

\subsection*{Excerpt Record 6 (Event 2 --- Closing Imperative Mobilizing "Nationalists of Bharat")}
\begin{itemize}[leftmargin=1.5em, itemsep=3pt, topsep=4pt]
  \item \textbf{Outlet:} Republic TV | \textbf{Date:} \texttt{2026-07-25}
  \item \textbf{Video title:} \textit{Rahul The Biggest Loser: Arnab Unravels The Political Agenda Behind CJP's Violent Protest}
  \item \textbf{Timestamp(s):} \texttt{[05:34--06:38]} | \textbf{Analyst:} Baani Solanki | \textbf{Language:} English (Roman script)
  \item \textbf{Exact transcript excerpt:}
    \begin{quote}\small\ttfamily [05:34] Dharmendra Pradhan has resigned when will Abhijit Dipke be held to account now [05:43] and if nationalist nationalist if you forgive and forget what they have done to our country this time, [05:52] then we will have to face it a second time. If you forget these pictures, then they will do it a third time. So, [06:01] learn your lesson and don't allow this CJP [ASR: BJP] to become a model... [06:10] Come together. India cannot forgive. India cannot forget. India will not forget or forgive all the [06:21] goons [ASR: goods] of Abhijit Dipke. Nationalists, mobilize yourself. Become more visible. [06:30] Become more vocal. Become more assertive. This nation is ours.\end{quote}
  \item \textbf{Feature(s):} $\boxtimes$ Naming $\boxtimes$ Presupposition $\boxtimes$ Agency $\boxtimes$ Antagonisation $\boxtimes$ Intensification / softening
  \item \textbf{Who is speaking? Is this the outlet's wording or a quotation?} Arnab Goswami's closing peroration addressed directly to \texttt{"Nationalists"}.
  \item \textbf{What does the wording imply, foreground, soften, or leave unstated?} Abandons journalistic reporting altogether for \textbf{political mobilization}: equates his ideological in-group (\texttt{"Nationalists"}) with \texttt{"India"} itself (\texttt{"India will not forget or forgive"}, \texttt{[06:10]}), and closes with four staccato imperatives (\texttt{"mobilize yourself. Become more visible. Become more vocal. Become more assertive"}) culminating in the exclusionary possessive \texttt{"This nation is ours"} (\texttt{[06:38]})---implying dissenting youth lie outside \texttt{"our"} nation!
  \item \textbf{What context or evidence supports or complicates this interpretation?} Reveals deep establishment anxiety that a decentralized, youth-led movement (\texttt{"don't allow this CJP to become a model"}, \texttt{[06:01]}) succeeded in toppling a Cabinet Minister, prompting a call for counter-street mobilization.
  \item \textbf{How does another outlet describe the same event or actor?} Contrast \texttt{"This nation is ours"} (Republic TV \texttt{[06:38]}, exclusionary partisan ownership) with \textbf{Mohak Mangal}'s title \texttt{"We the people of India..."} (\texttt{v-RikSqb0Vw}, inclusive constitutional citizenship) and \textbf{The Wire}'s Naseeruddin Shah (\texttt{pYaTVvHeGoM}, \texttt{[25:33--25:42]}: \texttt{"That house of parliament is ours. It belongs to the kids"}).
  \item \textbf{My interpretation, with reasons:} Comparing who \texttt{"ours"} refers to in Republic TV (\texttt{"Nationalists... This nation is ours"}) versus The Wire/Mohak Mangal (\texttt{"We the people... It belongs to the kids"}) captures the entire ideological fault line of contemporary Indian media discourse.
\end{itemize}
\vspace{4pt}\noindent\rule{\textwidth}{0.4pt}\vspace{4pt}

\section*{Common form: End-of-event summary}

\subsection*{End-of-Event Summary --- Event 1 (Ram Mandir Donation Controversy)}
\begin{itemize}[leftmargin=1.5em, itemsep=3pt, topsep=4pt]
  \item \textbf{The main way this outlet framed the event was:} As an isolated \texttt{Rs.\,80+ lakh} cash-counting theft by 8 \texttt{"low-level functionaries"} (\texttt{[03:33]}, \texttt{[07:14--07:26]}) swiftly punished under CM Yogi Adityanath's \texttt{"zero tolerance"} governance (\texttt{[01:12]}, \texttt{[03:22]}), which anti-Hindu opposition leaders (\texttt{Kejriwal}, \texttt{Akhilesh}, \texttt{Rahul}) hypocritically politicized---justifying putting the opposition \texttt{"in the dock"} (\texttt{[05:14--05:24]}) and condemning AAP's \texttt{"Taliban"} rhetoric against 79-year-old VHP leader Champat Rai (\texttt{[17:20--21:28]}).
  \item \textbf{The strongest linguistic evidence was (quote + timestamp):} \texttt{"there is no religious angle, there is no political angle, there is no government angle in this controversy... since they have politicized the matter I suggest the best response is to put them in the dock and ask them what their record is on Hindutva."} (\texttt{[04:08]} \& \texttt{[05:14--05:24]}).
  \item \textbf{The shared feature(s) most clearly illustrated were:} \textbf{Naming} (\texttt{low-level functionaries, some cashiers} vs \texttt{79-year-old prominent VHP leader}), \textbf{Agency} (agentless \texttt{"there has been theft"} vs active credit \texttt{"immediate action by the Yogi government"} at \texttt{[00:54--01:12]}), and \textbf{Intensification / Antagonisation} (\texttt{GENOCIDE OF KAR SEVAKS} at \texttt{[04:56]} and \texttt{"PART OF THE TALIBAN SHOOTING SQUAD"} at \texttt{[33:37]}).
  \item \textbf{An alternative interpretation of the same evidence is:} Articulated briefly by Arnab himself at \texttt{[39:47--40:16]} and by opposition panelists (\texttt{[14:03--15:23]}, \texttt{[27:34]}): that multi-year organized theft across 70+ CCTV instances and unreceipted silver gifts could not occur without senior Trust negligence or patronage, making quiet resignations without FIR naming an \texttt{"honorable exit"} rather than true \texttt{"zero tolerance"}.
  \item \textbf{What I cannot conclude from this transcript alone:} What the unpublished SIT report concluded regarding senior trustees' oversight, whether Raju Manwani's \texttt{200 kg} silver bricks or 2021 land-deal allegations were investigated, or whether custodial interrogation of the 8 clerks (\texttt{[50:13--50:22]}) implicated higher-ups.
\end{itemize}

\subsection*{End-of-Event Summary --- Event 2 (Students' Protest Regarding NEET Paper Leak)}
\begin{itemize}[leftmargin=1.5em, itemsep=3pt, topsep=4pt]
  \item \textbf{The main way this outlet framed the event was:} By \textbf{completely erasing NEET paper leaks and genuine students} (\texttt{"the students had evaporated"}, \texttt{[05:26]}), portraying the Jantar Mantar movement as a \texttt{"Deep State"} / Rahul Gandhi \texttt{"streetocracy"} conspiracy (\texttt{[00:23--00:51]}) of \texttt{"rabid, bloodthirsty mobs"} attacking RAF/police while soldiers died fighting terrorists in Anantnag (\texttt{[01:47]}, \texttt{[04:19--04:46]}), spinning Dharmendra Pradhan's resignation as a masterstroke that \texttt{"outplayed"} the opposition (\texttt{[02:06--02:56]}), and exhorting \texttt{"Nationalists"} never to forgive or forget (\texttt{[05:43--06:38]}).
  \item \textbf{The strongest linguistic evidence was (quote + timestamp):} \texttt{"our police, our RAF, our security personnel were vehemently attacked by rabid, bloodthirsty mobs, comes a resignation... these were not students, you were fooled... Nationalists, mobilize yourself... This nation is ours."} (\texttt{[01:37--01:47]}, \texttt{[04:58]} \& \texttt{[06:21--06:38]}).
  \item \textbf{The shared feature(s) most clearly illustrated were:} \textbf{Naming} (\texttt{Deep State Dipke}, \texttt{rabid, bloodthirsty mobs}, \texttt{Lynch mob}, \texttt{streetocracy}), \textbf{Agency} (deleting 100\% of police violence and hiding government capitulation inside agentless \texttt{"comes a resignation"} + 5x \texttt{"the government outplayed them"}), and \textbf{Antagonisation} (pitting \texttt{"Nationalists"} \& \texttt{"brave soldiers in Anantnag"} against protesters, closing with \texttt{"This nation is ours"}).
  \item \textbf{An alternative interpretation of the same evidence is:} That Dharmendra Pradhan resigned and the government dropped cases against Dipke and Saurabh Das (\texttt{[00:00--00:15]}) precisely because the movement represented massive, legitimate student anger over exam leaks---not a foreign \texttt{"Deep State"} mob (which MEA itself denied, \texttt{EM1GeZufO6E} \texttt{[00:40]}).
  \item \textbf{What I cannot conclude from this transcript alone:} Anything at all about the NEET examination crisis (never mentioned once!), student casualties from lathis/tear-gas/pellet guns on 20 July, or the terms of CJP's settlement beyond Pradhan's resignation and non-prosecution assurance (\texttt{[00:00--00:15]}).
\end{itemize}
\textit{Research reminder: interpret wording in context, distinguish allegations from established facts, and support claims with exact excerpts. Mixed or contradictory evidence is worth recording.}


\end{document}
